% !TeX program = xelatex
% Compile twice with XeLaTeX to resolve citations and cross-references.
\documentclass[a4paper,11pt]{article}
\usepackage[english]{babel}
\usepackage{fontspec}
\setmainfont{Latin Modern Roman}
\usepackage[
  top=26mm,bottom=25mm,left=25mm,right=25mm,
  headheight=14pt,headsep=6mm
]{geometry}
\usepackage{amsmath,array,booktabs,tabularx,caption,cite,fancyhdr,titlesec}
\usepackage{microtype}
\usepackage{hyperref}
\hypersetup{
  unicode,hidelinks,pdfdisplaydoctitle=true,
  pdftitle={A Review of Large Language Model Agents for Density Functional
    Theory Calculations with ORCA},
  pdfsubject={Task planning, tool use, execution feedback, and scientific
    validation in ORCA-based workflows},
  pdfkeywords={large language models, ORCA, agents, density functional theory,
    quantum chemistry, scientific workflows, reproducibility}
}
\titleformat{\section}{\large\bfseries\raggedright}{\thesection}{0.7em}{}
\titleformat{\subsection}{\normalsize\bfseries\raggedright}{\thesubsection}{0.7em}{}
\titlespacing*{\section}{0pt}{2.8ex plus .5ex}{1ex}
\titlespacing*{\subsection}{0pt}{1.8ex plus .4ex}{.5ex}
\pagestyle{fancy}
\fancyhf{}
\fancyhead[C]{\small Large Language Model Agents for DFT Calculations with ORCA}
\fancyfoot[C]{\small\thepage}
\renewcommand{\headrulewidth}{0.3pt}
\setlength{\emergencystretch}{2em}
\setlength{\parskip}{.15em}
\renewcommand{\arraystretch}{1.18}
\setlength{\tabcolsep}{5pt}
\captionsetup[table]{position=top,font=small,labelfont=bf,labelsep=quad,skip=6pt}
\newcolumntype{Y}{>{\raggedright\arraybackslash}X}
\newcolumntype{P}[1]{>{\raggedright\arraybackslash}p{#1}}
\newcommand{\doi}[1]{\href{https://doi.org/#1}{\nolinkurl{#1}}}
\newcommand{\ax}[1]{\href{https://arxiv.org/abs/#1}{arXiv:\,#1}}
\newcommand{\gitrepo}[2]{\href{https://github.com/#1}{#2}}
\newcommand{\snap}[2]{\href{https://github.com/#1/commit/#2}{\texttt{#2}}}
\urlstyle{same}
\raggedbottom

\begin{document}
\thispagestyle{plain}
\begin{center}
  {\LARGE\bfseries A Review of Large Language Model Agents for\\[.3ex]
   Density Functional Theory Calculations with ORCA\par}
\end{center}

\section{Introduction}
A molecular calculation can terminate normally and still leave the research
question unanswered.
For example, ORCA can compute the magnetic shielding constants of individual
nuclei in a target molecule.
Deriving a chemical shift also requires a reference value consistent with the
computational protocol~\cite{quntur}.
A quantum chemistry agent must identify the physical quantity requested in a
natural-language question, determine which calculations it requires, call the
solver, and use the results to assess whether the original plan was adequate.
This review examines which of these decisions a large language model (LLM)
makes and how software and researchers check them.

\section{Scope and Basis for Comparison}
The review covers 17 articles or preprints and one standalone benchmarking
software package.
It focuses on studies that used ORCA to perform density functional theory
(DFT) calculations.
Studies of ORCA input generation and multireference wavefunction calculations
provide related evidence.
For systems that support several programs, we classify each example by the
solver and method actually used.
We examine who constructs the plan, what feedback changes the next action,
and what evidence supports the reported conclusion.
In reporting controlled comparisons, we retain each study's unit of evaluation,
baseline, and assessment criteria.
We compare numerical results only across matched conditions within a study.

\section{How ORCA Agents Plan and Execute Calculations}

\subsection{From Input Generation to Solver Execution}
Jacobs and Pollice established a baseline for generating ORCA input files.
On 294 test prompts derived from actual ORCA inputs,
5.872\% of the files generated by the GPT-4o baseline passed an execution check.
With fine-tuning on synthetic data and chain-of-thought prompting,
the proportion rose to 18.027\%.
For this check, the authors appended the same water coordinates to each input
and required ORCA to return an exit code of zero within five seconds.
The test measures whether a model can translate a request into input the solver
will accept.
Verifying the geometry of the original molecule, dependencies among calculations,
and the target property requires evaluating the full task~\cite{input}.

ArIA integrates input generation into an interactive workflow.
It interprets a user's instructions, runs ORCA, produces figures and reports,
and accepts follow-up questions about the results.
The authors provide separate resources for model development and deployment.
During development, the authors used ORCA trial runs together with validation
and feedback modules to check candidate inputs.
In the deployed system, the trained model adapters are integrated into a
LangGraph workflow~\cite{aria}.
These resources support separate evaluations of input quality and of how well
follow-up questions help users interpret a particular calculation.
Each evaluation requires records of the corresponding inputs, outputs, and
interactions.

\subsection{Plan Representation and Validation before Submission in Q-Planner}
Q-Planner represents a user's request as a directed acyclic graph.
The model identifies the compounds, selects a computational protocol, and
specifies the inputs and outputs of each node.
Once the plan passes structural validation, an executor schedules tool calls
in dependency order, invokes ORCA through the ORCA Python Interface, and
passes the results to downstream nodes.
The plan can specify anticipated errors and the corresponding parameter changes.
The executor can run independent tasks in parallel and apply these rules to
known failures~\cite{qplanner}.

The aqueous $\mathrm{p}K_\mathrm{a}$ calculation for chlorofluoroacetic acid
illustrates this division of responsibility.
The model omitted geometry optimization from its first plan.
The validator detected the omission before any ORCA job was submitted.
A second planning call restored the missing step, allowing the free-energy
calculations that depended on optimized structures to proceed~\cite{qplanner}.
A separate experiment tested batch planning for single-point calculations on
sets of 10 to 100 molecules.
All 30 plans generated by expanding templates contained every required node,
compared with 28 of 30 plans generated node by node.
These counts measure completeness at the level of the plan.
Assessing the reliability of the energy for any individual molecule requires
the executed job and its output.

\subsection{Task Delegation and Execution Feedback in El Agente Q and Quntur}
El Agente Q uses hierarchical delegation.
A computational chemist agent breaks the question into steps.
The agent responsible for DFT then delegates input preparation and cluster
submission to separate specialists.
After calling their tools, these specialists report results or failures to the
supervising agent, so that subsequent steps can reuse existing files and the
plan can be revised when a job fails.
The authors released calculation records for six exercises and two case studies.
The exercises specify the functional and basis set in advance.
They therefore test chiefly whether the agent can translate a prescribed method
into a coherent sequence of ORCA jobs~\cite{agente}.

Quntur extends this feedback loop to method selection and information retrieval.
During input preparation, the coordinating agent identifies whether it needs
information about program syntax or a scientific justification for a choice.
It can then consult the ORCA manual, search the literature, or delegate to a
domain specialist.
For long ORCA outputs, a tool first indexes section names and line ranges.
The agent uses this index to locate a table, inspect its header and sample
rows, and extract the fields needed for analysis~\cite{quntur}.
A single output file can contain several electronic states or properties.
Section and line references allow researchers to check that the extracted data
correspond to the requested quantity.

\section{How ORCA Output Informs an Agent's Scientific Judgments}
Quntur's analysis of 4-(dimethylamino)benzonitrile (DMABN) illustrates how
computational results can change an interpretation.
The agent used optimized geometries and natural transition orbitals to determine
whether two low-lying singlet states had locally excited or charge-transfer
character.
When the change in state character differed from the published account, the
agent searched the literature again.
It traced the discrepancy to the conditions being compared.
The calculation treated the molecule in the gas phase, while the literature
comparison concerned the molecule in solution~\cite{quntur}.
This finding qualified the interpretation and identified the environmental
condition that a follow-up calculation would need to control.

In a comparison covering 17 tasks, each repeated independently five times,
Quntur and Claude Code with ORCA-specific skills both used Claude Opus 4.5.
Experts scored planning, geometry preparation, input generation and execution,
and post-processing and interpretation.
Mean scores were 90.5 out of 100 for Quntur and 61.7 for Claude Code.
On the nuclear magnetic resonance and magnetic properties task, the baseline
omitted one or both reference calculations in four of five runs~\cite{quntur}.
This example links the score to a verifiable decision about whether to schedule
the reference calculation required for the requested property.
The overall mean reflects the complete configuration.
Estimating the contribution of an individual module would require an ablation
study.

Lee and Rondinelli used a different evaluation design in their study of
multireference wavefunction calculations with ORCA.
The agent selected active spaces, submitted jobs, handled errors, and matched
electronic states to the targets.
With expert guidance on these decisions, the number of matched reference
transitions increased from 139 to 246 out of 558.
For the 109 transitions matched by both configurations, the mean absolute
error changed from 0.242 to 0.268~eV~\cite{mr}.
This finding illustrates why evaluations should report both the number of
targets with usable results and the numerical error on the same set of
matched targets.

\section{Lessons for Agent Design from Other Computational Platforms}
Molecular workflows with multiple back ends illustrate other ways to assign
tasks to tools.
ChemGraph decomposes tasks and executes them through the Atomic Simulation
Environment (ASE). Its small-molecule geometry optimization benchmarks use
NWChem reference calculations.
Aitomia selects methods and executes scripts through MLatom. Its infrared
spectrum example for hexanol uses AIQM2~\cite{chemgraph,aitomia}.
Cloudra checks saddle points located with Gaussian by verifying the endpoints
of the reaction path.
In VirtualLab\_CC, the principal investigator approves progression to the next
stage. Its published protocol assigns DFT conformer screening to Gaussian
and DLPNO-CCSD(T) reference calculations to ORCA~\cite{cloudra,virtuallab}.
These cases illustrate how tool selection, independent checks, and human
review enter a computational task.

Studies of periodic systems address when an agent should intervene again.
VASPilot revises inputs in response to job logs, and LatticeMind checks data
dependencies between steps before submission~\cite{vaspilot,lattice}.
In an AutoDFT comparison using the same model, closed-loop feedback increased
the number of completed VASP tasks from 28 to 32 out of 34.
Mean computational cost rose from 22.1 to 44.0 H100 GPU-hours per
task~\cite{autodft}.
CatMaster revises the candidate pool in light of DFT checks, and MOFGen organizes
candidate generation, screening, and verification into separate
stages~\cite{catmaster,mofgen}.
DFT-Agent Bench examines solver execution records, and CompChemBench assesses
the ability to audit a workflow.
Work on functional discovery evaluates the predictive performance of
agent-proposed functionals on held-out data~\cite{dftbench,compbench,xc}.

\section{Discussion and Limitations}
The ORCA studies suggest that the value of feedback depends on the decision
it can correct.
A trial run can expose input syntax errors, and plan validation can detect
a missing upstream geometry optimization.
Recognizing a missing reference state or a mismatch between gas-phase and
solution conditions requires relating the actual output to the user's question.
For tasks with known dependencies, the Q-Planner comparison supports expanding
templates during planning.
When the interpretation depends on intermediate results, the Quntur case
illustrates when it is useful to consult the literature again~\cite{qplanner,quntur}.

The studies differ in tasks, expert scoring rubrics, and computational budgets.
A high score on one task set therefore provides limited evidence of an
advantage on another.
A more informative comparison would hold the model, tools, and budget fixed,
record the agent's key decisions, and report task coverage, property errors
on a common set of cases, repeated calculations, and resource consumption.
Such comparisons would help determine whether gains arise from more complete
plans, better-targeted corrections, or more opportunities to run
calculations~\cite{quntur,mr,autodft}.

\section{Future Directions for the Project}

When an ORCA job finishes, researchers using an agent for molecular calculations
often face two successive questions.
Which follow-up questions can the existing results answer?
If further calculations are needed, which candidates should receive priority
given the available computing time?
Archived calculations on ethanol and \emph{n}-hexane motivate the two studies
proposed below.
The blinded test sets and controlled comparisons for these studies have yet
to be established.

\subsection{Reusing Calculation Outputs to Answer Follow-up Questions}

The conclusions that a given set of outputs can support depend on the physical
quantity of interest.
In Quntur's nuclear magnetic resonance benchmark, omitting a reference calculation
compromised the derivation of chemical shifts even when the shielding constants
of the target molecule were available~\cite{quntur}.
Lo Scudo tested a check on the completeness of thermochemical records for
individual species. When data for a required species were missing, the check
prevented the available values from being combined into a reaction
quantity~\cite{tracestate}.
The proposed study would test whether, given the same raw files and a different
question, an agent can determine which conclusions the evidence supports
and identify the conditions or data needed to draw a further inference.

The archived ethanol calculations illustrate the distinction between two
closely related questions.
B3LYP single-point electronic energies were evaluated at a common molecular
geometry in the gas phase and in water represented by the conductor-like
polarizable continuum model (CPCM).
For $\Delta E=E_{\mathrm{water}}-E_{\mathrm{gas}}$, the def2-SVP and
def2-TZVP basis sets gave $-16.61$ and $-18.28$~kJ\,mol$^{-1}$, respectively.
At this geometry, both basis sets therefore predict a lower B3LYP electronic
energy in implicit water than in the gas phase.
The two estimates of $\Delta E$ differ by approximately
$1.67$~kJ\,mol$^{-1}$~\cite{projectarchive}.
A question about the standard free energy of hydration requires additional
specification of the temperature, standard states, and protocol for calculating
the solvation free energy, followed by a check of which contributions are
present in the existing outputs.

The proposed blinded evaluation would use pairs of questions about molecules
held out during development.
Both questions in a pair would refer to the same fixed set of ORCA outputs.
They would differ only in the direction of subtraction, the reference
conditions, or the target quantity.
Independent annotators would supply answers that can be read from the raw files
or obtained by post-processing and identify questions requiring further
information or a new calculation.
With the model, tools, and budget held constant, the existing system would be
compared with a procedure that explicitly checks the target quantity, the
calculation conditions, and the source of each numerical value.
The primary outcome would be the proportion of pairs in which the agent
correctly assesses the evidence and interprets the physics for each question.
Additional outcomes would include the rate of numerical answers given on
insufficient evidence, the rate of unanswered questions for which the evidence
is sufficient, and the number of avoidable ORCA runs.

\subsection{Adaptive Conformer Selection under a Limited Computational Budget}

In multilevel conformer screening, an inexpensive method ranks a large set of
candidates before a subset is refined with DFT.
CENSO applies this approach to conformer ensembles and free-energy calculations
for flexible molecules~\cite{censo}.
VirtualLab\_CC reports the computing time spent on conformer screening and
subsequent refinement.
Controlled comparisons in AutoDFT also show that allowing an agent to revise
its plan in response to intermediate results can change the computational
cost~\cite{virtuallab,autodft}.
The proposed study would test whether adapting the selection of the next
candidate to the energies already obtained improves the probability of
identifying low-energy conformers for a fixed number of high-level
single-point calculations.

Two archived batches of gas-phase $r^2$SCAN-3c single-point calculations on
\emph{n}-hexane provide electronic-energy rankings within sets of four and six
fixed geometries, respectively~\cite{projectarchive}.
For the proposed evaluation, each additional small molecule would have a
fixed set of more than four candidate geometries.
Single-point energies for all candidates would be computed at a common
reference level of theory and withheld from the selection procedures.
Each procedure would receive the rankings from two low-cost methods and would
be allowed four queries of the reference energies.
The first baseline would select the four candidates with the lowest energies
according to the primary screening method.
The second would also seek coverage across conformer families among low-energy
candidates, using predefined dihedral-angle classes.
The adaptive procedure would begin with the lowest-energy candidate under the
primary screening method.
It would then choose each subsequent candidate using the reference energies
already returned, disagreements between the two rankings, and the conformer
families still unrepresented.

Once the reference energies are revealed, the hit rate would be calculated as
the fraction of molecules for which the procedure selected a conformer with
the lowest reference energy in the candidate set.
For a fixed candidate set $S$ and a final choice $\hat c$, the energy difference
$E_{\mathrm{ref}}(\hat c)-\min_{c\in S}E_{\mathrm{ref}}(c)$
measures the energetic cost of missing the lowest-energy conformer.
Reporting actual CPU core-hours alongside model token usage and grouping
molecules by the agreement between the two low-cost rankings would help
identify the molecules for which adaptive selection is beneficial.
The resulting conclusions would apply to electronic-energy rankings at the
specified geometries and reference level of theory.

\section{Conclusions}

Studies of ORCA workflows have demonstrated two approaches to agent design.
When dependencies are known, the model can produce a plan for validation
before submission.
A deterministic executor can then expand batch tasks and handle predefined
error conditions.
When intermediate results change the scientific question, the agent must
retain the raw outputs and identify missing references or calculation
conditions before deciding whether further calculations are needed and how
to interpret the results.
Comparative evaluations should record where a plan was corrected, which
decisions changed subsequent actions, and how those decisions affected both
the scientific conclusions and the computational cost~\cite{qplanner,agente,quntur}.

\begingroup
\small
\begin{thebibliography}{99}
\raggedright
\setlength{\itemsep}{.35em plus .1em}
\bibitem{quntur}
  PÉREZ-SÁNCHEZ J B, ZOU Y, CAMPOS-GONZALEZ-ANGULO J A, et al.
  El Agente Quntur: A research collaborator agent for quantum chemistry [J/OL].
  \emph{Chem}, published online 24 September 2026, article 103259.
  DOI: \doi{10.1016/j.chempr.2026.103259}.
\bibitem{input}
  JACOBS P F, POLLICE R.
  Developing large language models for quantum chemistry simulation input generation [J].
  \emph{Digital Discovery}, 2025, 4(3): 762--775.
  DOI: \doi{10.1039/D4DD00366G}.
\bibitem{aria}
  CHANMUNGKALAKUL S, VAN DER REE M, GIUNTOLI A, et al.
  Molecular Simulations Assisted by an Artificial Intelligence Agent (ArIA) [EB/OL].
  \emph{ChemRxiv}, v2, 2026-05-08.
  DOI: \doi{10.26434/chemrxiv.15002344/v2}.
\bibitem{qplanner}
  ZHANG R, GAO W, KURASHIGE Y.
  Q-Planner: A Harness System for Quantum Chemistry Agents [J/OL].
  \emph{ACS Omega}, published online 10 November 2026.
  DOI: \doi{10.1021/acsomega.6c05744}.
\bibitem{agente}
  ZOU Y, CHENG A H, ALDOSSARY A, et al.
  El Agente: An autonomous agent for quantum chemistry [J].
  \emph{Matter}, 2025, 8(7): 102263.
  DOI: \doi{10.1016/j.matt.2025.102263}.
\bibitem{mr}
  LEE V C, RONDINELLI J M.
  Can Autonomous LLM Agents Execute Multireference Quantum Chemistry Calculations?
  [EB/OL].
  \emph{arXiv}, v1, 2026-09-11.
  \ax{2609.13357v1}.
\bibitem{chemgraph}
  PHAM T D, TANIKANTI A, KEÇELİ M.
  ChemGraph as an agentic framework for computational chemistry workflows [J].
  \emph{Communications Chemistry}, 2026, 9(1): 33.
  DOI: \doi{10.1038/s42004-025-01776-9}.
\bibitem{aitomia}
  HU J, NAWAZ H, HOU Y-F, et al.
  Aitomia: An Agentic Framework for AI-Driven Atomistic and Quantum Chemical Simulations
  [J].
  \emph{Journal of Chemical Theory and Computation}, 2026, 22(13): 6557--6571.
  DOI: \doi{10.1021/acs.jctc.6c00591}.
\bibitem{cloudra}
  YANG C, ZHOU Q, LI Y.
  Cloudra: Geometric Reasoning and Scientific Validation for Autonomous Reaction Mechanism
  Exploration [EB/OL].
  \emph{arXiv}, v1, 2026-11-02.
  \ax{2611.02909v1}.
\bibitem{virtuallab}
  KIENINGER M, SEGOVIA M, VEGA-TEIJIDO M, et al.
  VirtualLab\_CC: An LLM-augmented virtual laboratory for automated computational
  chemistry research [J].
  \emph{Artificial Intelligence Chemistry}, 2026, 4(2): 100129.
  DOI: \doi{10.1016/j.aichem.2026.100129}.
\bibitem{vaspilot}
  LIU J, ZHU T, YE C, et al.
  VASPilot: MCP-facilitated multi-agent intelligence for autonomous VASP simulations [J].
  \emph{Chinese Physics B}, 2025, 34(11): 117106.
  DOI: \doi{10.1088/1674-1056/ae0681}.
\bibitem{lattice}
  VEGH N M, ZHU Y, PRENTKI R J, et al.
  LatticeMind: A Validation-Gated Agentic Framework for Density Functional Theory
  Workflows [EB/OL].
  \emph{arXiv}, v1, 2026-07-19.
  \ax{2610.26869v1}.
\bibitem{autodft}
  YANG P, ZHANG Z, LI Y, et al.
  AutoDFT: A Closed-Loop Multi-Agent Framework for Autonomous DFT Calculations [EB/OL].
  \emph{arXiv}, v2, 2026-06-04.
  \ax{2605.26179v2}.
  \href{https://phyang.top/}{Accepted at KDD 2027}.
\bibitem{catmaster}
  CHEN H, QIU J, TEW Y S, et al.
  Autonomous computational catalysis through an agentic research system [EB/OL].
  \emph{arXiv}, v4, 2026-06-05.
  \ax{2601.13508v4}.
\bibitem{mofgen}
  INIZAN T J, YANG S, KAPLAN A, et al.
  System of agentic AI for the discovery of metal--organic frameworks [J].
  \emph{Nature Computational Science}, 2026, 6(11): 1223--1234.
  DOI: \doi{10.1038/s43588-026-01058-x}.
\bibitem{dftbench}
  LIU S, YANG J, WAN W, et al.
  DFT-Agent Bench: Scaffold Effects and Execution Provenance in Autonomous DFT [EB/OL].
  \emph{arXiv}, v1, 2026-11-07.
  \ax{2611.10629v1}.
\bibitem{compbench}
  WU C.
  CompChemBench: A Benchmark for AI Agents on Computational Chemistry Tasks [CP/OL].
  v1.1. \emph{GitHub}, 2026 [2026-11-22].
  \href{https://github.com/Chenghao-Wu/compchem-bench/tree/772d3f3e1a93}{Commit 772d3f3e1a93}.
\bibitem{xc}
  CHANDRASEKHAR A, BARATI FARIMANI A.
  Agent-Guided Discovery of Exchange-Correlation Functionals with Physics-based Validation
  [EB/OL].
  \emph{arXiv}, v1, 2026-11-09.
  \ax{2611.13329v1}.
\bibitem{tracestate}
  LO SCUDO F.
  Trace-State Validation for Agentic Computational Chemistry Workflows [EB/OL].
  \emph{ChemRxiv}, v1, 2026-07-20.
  DOI: \doi{10.26434/chemrxiv.15006307/v1}.
\bibitem{projectarchive}
  BG6022-V4 Project Team.
  Calculation records and report verification for conformers of \emph{n}-hexane
  and environmental effects on ethanol [R/OL, in Chinese].
  \emph{GitHub}, 2026.
  \href{https://github.com/az87988799/BG6022-v4/blob/1e8795b0b10deb8f288eccb9d95befb742958ff9/FinalReport-LaTeX/sections/06a-calculation-results.tex}
    {Calculation archive} and
  \href{https://github.com/az87988799/BG6022-v4/blob/1eefccdb742ed81692900f9660a8db950ae1f273/docs/acceptance/live-failure-repair-20261021/repair-report.md}
    {verification record for the ethanol calculations}.
\bibitem{censo}
  GRIMME S, BOHLE F, HANSEN A, et al.
  Efficient Quantum Chemical Calculation of Structure Ensembles and Free Energies for
  Nonrigid Molecules [J].
  \emph{The Journal of Physical Chemistry A}, 2021, 125(19): 4039--4054.
  DOI: \doi{10.1021/acs.jpca.1c00971}.
\end{thebibliography}
\endgroup

\clearpage
\appendix
\setcounter{table}{0}
\renewcommand{\thetable}{\Alph{section}.\arabic{table}}
\renewcommand{\theHtable}{appendix.\Alph{section}.\arabic{table}}
\section{Public Software and Reproducibility Resources}
\label{app:repos}
The resources released for ORCA-based agents document different aspects of
their workflows.
The Q-Planner code snapshot includes skill specifications, the plan validator,
the task scheduler, and the service that exposes ORCA tools.
These components can be checked against the plan files to trace the inputs
and outputs of individual nodes~\cite{qplanner}.
The \href{https://doi.org/10.5683/SP3/JU2BQK}
  {archive of eight El Agente Q tasks}
and the \href{https://doi.org/10.5683/SP3/RDSOEA}
  {evaluation archive for Quntur and its baseline}
contain the files produced during execution, allowing readers to trace the
reported values back to the agents' decisions.
The Quntur authors report that the source code for the agent is not publicly
available~\cite{agente,quntur}.
The \href{https://doi.org/10.34894/RNPTDS}
  {ArIA deployment and training materials}
document its input-generation and feedback modules.
Jacobs and Pollice provide an
\href{https://git.lwp.rug.nl/pollice-research-group/ORCAInputFileSynthesis/-/blob/f60237db580d44849f30179deed090317a4a768a/Classes/Helpers/ORCARunner.py}
  {ORCA test runner}
and an \href{https://doi.org/10.34894/WNRHA4}{archived dataset}~\cite{aria,input}.

Table~\ref{tab:repos} lists 12 public GitHub repositories.
For Q-Planner, the snapshot corresponds to the publication tag
\texttt{qplanner-acsomega-r2.2}.
Each commit identifier links to the version examined for this review.

\begin{table}[htbp]
\centering
\caption{Public code and data repositories included in the comparison
  (accessed 22 November 2026).}
\label{tab:repos}
\small
\begin{tabularx}{\linewidth}{@{}P{3.45cm}YP{2.95cm}@{}}
\toprule
Software or dataset
  & Agent or evaluation components available for inspection
  & Commit\\
\midrule
\gitrepo{zenqiaz/Q-planner}{Q-Planner}
  & Plan graphs, the validator, and logs of ORCA tool calls~\cite{qplanner}
  & \snap{zenqiaz/Q-planner}{570409b96fa8}\\
\gitrepo{argonne-lcf/ChemGraph}{ChemGraph}
  & Task decomposition, ASE tool calls, and aggregation of results~\cite{chemgraph}
  & \snap{argonne-lcf/ChemGraph}{dcb129ef6269}\\
\gitrepo{dralgroup/aitomia}{Aitomia}
  & Method selection, script generation, and execution in MLatom~\cite{aitomia}
  & \snap{dralgroup/aitomia}{e9d62acafa53}\\
\gitrepo{yyywaa/Cloudra-mechanism-agent}{Cloudra}
  & Planning of mechanistic studies and verification of reaction-path endpoints~\cite{cloudra}
  & \snap{yyywaa/Cloudra-mechanism-agent}{d76d5cca194c}\\
\gitrepo{CCBG-Lab/VirtualLab_CC}{VirtualLab\_CC}
  & Research roles, human review, and the
    \href{https://github.com/CCBG-Lab/VirtualLab_CC/blob/1c1a58cc8df6/framework/methods/METH_002_criegee_terpene_pes.md}
    {computational protocol}~\cite{virtuallab}
  & \snap{CCBG-Lab/VirtualLab_CC}{1c1a58cc8df6}\\
\gitrepo{JiaxuanLiu-Arsko/VASPilot}{VASPilot}
  & Job monitoring and error recovery using run logs~\cite{vaspilot}
  & \snap{JiaxuanLiu-Arsko/VASPilot}{2d9b2dfaf6e0}\\
\gitrepo{penghui-yang/AutoDFT}{AutoDFT}
  & Closed-loop execution and the open-loop control~\cite{autodft}
  & \snap{penghui-yang/AutoDFT}{4c8e16ff59b7}\\
\gitrepo{nathaniel-vegh/latticemind-data-package}{LatticeMind data package}
  & Input dependency checks and ablation data~\cite{lattice}
  & \snap{nathaniel-vegh/latticemind-data-package}{2a85c9ef5c2b}\\
\gitrepo{q734738781/CatMaster}{CatMaster}
  & Tools for screening and verifying catalyst candidates~\cite{catmaster}
  & \snap{q734738781/CatMaster}{06b856814f98}\\
\gitrepo{tjaffrel/mofgen}{MOFGen}
  & Candidate generation and multistage screening~\cite{mofgen}
  & \snap{tjaffrel/mofgen}{841c93995a4c}\\
\gitrepo{Chenghao-Wu/compchem-bench}{CompChemBench}
  & Task execution and scoring of scientific audits~\cite{compbench}
  & \snap{Chenghao-Wu/compchem-bench}{772d3f3e1a93}\\
\gitrepo{BaratiLab/DFT_functional_agentic_searcher}{Exchange--correlation functional search}
  & Proposing candidate functionals and testing them on independent data~\cite{xc}
  & \snap{BaratiLab/DFT_functional_agentic_searcher}{741bd205c5a9}\\
\bottomrule
\end{tabularx}
\end{table}

\end{document}
